\section{A selected physical roof height for an actual clearance witness}
\label{sec:lc-roof}

We now use the length-sensitive source estimate at the same physical
roof as the observation.  Keep the complete cap
$\mathcal C_{m,R}$, its strict regular convention, and the original
normalization of \eqref{eq:rbv-majorant}.  Thus $a=5/6$, every original
exact event is a positive restriction of this cap source, and the
unconditioned collision measure is divided by $c_*$ exactly once.
For $m\ge1$, $0\le j<m$ and $0<\delta\le\delta_*=3/200$, define
\begin{equation}\label{eq:lc-selected-roof}
 E_{j,\delta,R}=\{x:N_{\delta,R}(T_R^jx)\ge1\},\qquad
 \mu_{m,R;j,\delta}
   =(L_{m,R})_*\bigl(\one_{\mathcal C_{m,R}}
                        \one_{E_{j,\delta,R}}\,d\nu/c_*\bigr).
\end{equation}
The index $j$ denotes an actual flight between collisions $j$ and $j+1$.
No collision is inserted at a closest point.  The measure is a positive
majorant on the original roof line; it is not a new conditional law.

\subsection{Turning constants for the enlarged physical source class}

The enlarged class, cutoff $\widehat k_0$, short length $\ell$, and
constants in Section~\ref{sec:lc-length} are used throughout this
section.  To avoid reusing a constant for a smaller curve class, set
\begin{align}
 B_2&=\widehat B_{\rm cv}(1+U^2)^{3/2},\notag\\
 a_0&=2k_-t_*,\qquad L_h=3\sqrt{1+U^2}/k_-,\notag\\
 c_H&=\sin(\widehat k_0^{-2}),\qquad
 c_a=\sqrt{1-a^2}=\sqrt{11}/6,\notag\\
 D_B&=t_*^{-1}+2/(t_*a_0)
              +(B_2/a_0)(1+(2k_-)^{-1})+B_2,\notag\\
 V_H&=\max\{-2\log c_H,
             \log(\pi/2)+2\log(1+\widehat k_0^{-1})\},
                                                   \label{eq:lc-turn-one}\\
 D_{\rm var}&=\max\{(U/c_H+D_B)L_0^{2/3},
                         (2\pi U+D_B)L_h^{2/3}\},\notag\\
 V_J&=D_{\rm var}\Gamma^{1/3}L_0^{1/3}/(1-\sigma^{1/3}),\notag\\
 V_\rho&=V_H+(a_\circ^{-1}+b_*a_\circ^{-2}+K_0/2)\ell,\notag\\
 V_t&=-2\log c_a+B_2L_0/k_-,\notag\\
 V_1&=L_0(k_+b_*+1)/(k_-a_\circ^2),\notag\\
 V_A^{\rm lc}&=\max\{V_1,V_\rho+V_J+V_t\},\notag\\
 T_{\rm lc}&=2\pi+V_A^{\rm lc}.
                                                   \label{eq:lc-turn-two}
\end{align}
Here $L_0$, $\Gamma$, $U$ and $\sigma$ have the meanings fixed above,
$a_\circ=\sqrt{1827}/200$, $b_*=97/200$, and $K_0$ is the exact
first-image derivative bound in \eqref{eq:ag-geometry-constants}.
All these constants are fixed before a collision, radius, label, level
or roof supremum.

\begin{lemma}[Turning on the actual selected interval covers]
\label{lem:lc-turning}
On each overlap of the two physical half-source interval partitions
of Theorem~\ref{thm:lc-local-half}, use the original cap and set
$P(\vartheta)=(-p_0(\vartheta),p_m(\vartheta))$.  The complete
punctured turning integral on that overlap obeys
\begin{equation}\label{eq:lc-turning}
                  \int_{\{P\ne0\}}|d\arg P|\le T_{\rm lc}.
\end{equation}
This includes the selected covers of that theorem and the strict
unvisited trace at $u_0=0$.  Every retained Borel observation may be
inserted as a selector in this positive integral.
\end{lemma}
\begin{proof}
We reproduce the variation argument for the enlarged class.  On an
actual homogeneous image segment of arclength $r$ in a signed strip
$H_k$ one has
\[
 r\le L_h k^{-3},\qquad c\ge (2/\pi)(k+1)^{-2}.
\]
The exact one-step logarithmic derivative \eqref{eq:uc-logjac}, with
\eqref{eq:lc-turn-one}, has absolute value at most $U/c+D_B$ per
unit image arclength.  Its absolute integral is at most
$(2\pi U+D_B)k^2r\le(2\pi U+D_B)L_h^{2/3}r^{1/3}$.
On the central strip the same integral is at most
$(U/c_H+D_B)L_0^{2/3}r^{1/3}$.  Each unstable graph has total
arclength at most $L_0$, because $V\ge k_-$ and its angle ranges
in $(-\pi/2,\pi/2)$.  Adapted expansion gives the intermediate
length bound $L_i\le\Gamma\sigma^{l-i}L_l$.  Summing the absolute
integrals of the actual one-step logarithms therefore bounds the
variation of a full homogeneous arclength Jacobian by $V_J$.
This is an estimate of total variation; it is not inferred merely
from a bound on differences of endpoint values.

For a half-ray with $l\ge2$ true collisions, its exact angular
first-image density is
$\rho_C=c_1/(a^{\rm foot}\sqrt{1+V_1^2})$.
The monotonicity of $\varphi_1$ and the strip bounds give variation
of $\log c_1$ at most $V_H$.  The exact first-flight identities
\eqref{eq:ag-a-derivative} give
$|(a^{\rm foot})_s|\le1+b_*/a_\circ$ and
$|V_{1,s}|\le K_0$.  Thus
$\operatorname{Var}\log\rho_C\le V_\rho$.  This step retains the
first grazing factor.  The arclength chain rule is exactly
\[
 A_l^p=|\partial_\vartheta p_l|
   =\frac{c_lV_l}{\sqrt{1+V_l^2}}
                   \frac{J_C(T_R^{l-1})}{\rho_C}.
\]
The terminal cap is one interval, since the physical endpoint
momentum is strictly monotone.  It lies inside the central strip
for every $\widehat k_0\ge2$.  On it $c_l\ge c_a$, and geometric
curvature gives $|V_{l,s}|\le B_2$.  Consequently the terminal
factor has logarithmic variation at most $V_t$.  The preceding
three bounds give
$\operatorname{Var}\log A_l^p\le V_\rho+V_J+V_t$.

For $l=1$ the actual incoming endpoint derivative is
$A_1^p=\kappa a^{\rm foot}+c$.  With signed impact
$b=\varsigma(R+u_0)$, direct differentiation gives
\[
 \partial_s(\kappa a^{\rm foot}+c)
                 =c(\kappa b-p)/a^{\rm foot}.
\]
Its logarithmic derivative has magnitude at most
$(k_+b_*+1)/(k_-a_\circ^2)$, hence variation at most $V_1$.
Physical reversal proves the same statement for the backward half.
Refining a connected branch by any of the interval covers in
Section~\ref{sec:lc-length} only decreases these variations.

On an overlap, both coordinate derivatives of $P$ have nonzero
fixed signs by \eqref{eq:bf-signed-derivatives}.  Its tangent remains
in one open quadrant, and $P$ has at most one zero.  Its tangent
angle, as a function of the logarithm of the absolute slope ratio,
has derivative bounded by $1/2$.  Since each half logarithmic
variation is at most $V_A^{\rm lc}$, the tangent turning is at
most $V_A^{\rm lc}$.  Lemma~\ref{lem:ab-radial-turning} gives
\eqref{eq:lc-turning}, with $2\pi$ allowing a single zero.  This
punctured estimate omits a polar-angle jump at zero; it does not
assign the separate incoming critical fan a zero value.  Borel
selectors multiply the bounded positive integral and are never
differentiated.  Compact regular exhaustion proves all open-interval
and one-sided trace statements.
\end{proof}

\subsection{All outgoing boundary charges of the selected source}

For a regular physical root write $u_x$ for the absolute inverse
endpoint-momentum Jacobian, and $h_x=R^2e^{\mathsf T}G^{-1}e$.
The exact identity of Corollary~\ref{cor:bf-wronskian} gives
\begin{equation}\label{eq:lc-boundary-density}
 A=\frac{u_x}{4\pi c_*h_x}\le C_Au_x,\qquad
 C_A=\frac{\Lambda_\circ}{4\pi c_*R_-^2}.
\end{equation}
Set
\begin{align}
 C_{\rm foot}&=607/90,&
 K_{\rm bd}&=C_AT_{\rm lc}C_{\rm foot}E_{\rm lc},\notag\\
 C_{\rm side}&=\frac{K_{\rm bd}(2K_{\rm loc}+H_{\rm lc})}
                         {1-\sigma},&
 C_{\delta}&=2K_{\rm bd}H_{\rm lc},\notag\\
 M_{\rm cap}&=\max\{A_{\rm src}G_\ell L_{\rm c},
                                      J_0L_{\rm c}/R_-\},&
 C_{\rm cap,loc}&=4C_AM_{\rm cap}/a.
                                                   \label{eq:lc-boundary-constants}
\end{align}
All marked sums in these constants include both impact signs.

\begin{lemma}[Selected old clearance boundaries and new level boundaries]
\label{lem:lc-selected-boundary}
The total regular boundary charge of the selected source
\eqref{eq:lc-selected-roof} on all original unvisited tangencies is
at most
\begin{equation}\label{eq:lc-old-boundary}
 C_{\rm side}\{\delta+\theta_{\rm r}^{j}
                             +\theta_{\rm r}^{m-1-j}\}.
\end{equation}
The additional regular boundary at clearance level $\delta$ on
flight $j$ has charge at most
\begin{equation}\label{eq:lc-new-boundary}
 C_{\delta}\{\theta_{\rm r}^{j}
                             +\theta_{\rm r}^{m-1-j}\}.
\end{equation}
These are complete positive boundary sums before any original
exact-label or roof restriction.
\end{lemma}
\begin{proof}
Consider an old unvisited tangency on flight $k$.  The two actual
half-lengths from its unreflected foot are $k+1$ and $m-k$.
If $k<j$, the selected flight is in the forward half.  In
\eqref{eq:lc-half-index} the precise indices are
\[
 h=j-k-1,\qquad n=m-j-1,\qquad l=m-k.
\]
Both $h$ and $n$ are nonnegative.  The backward half is unselected.
The interval-overlay inequality \eqref{eq:mb-overlay}, the
selected capacity \eqref{eq:lc-local-half}, and the individual
bound \eqref{eq:lc-individual} give
\[
 \sum_{\text{actual selected overlaps}}\sup u_x
 \le C_{\rm foot}E_{\rm lc}
 \{K_{\rm loc}\sigma^k(\delta+\theta_{\rm r}^{m-j-1})
                             +H_{\rm lc}\sigma^{m-k-1}\}.
\]
Indeed the exact determinant bounds each root by
$C_{\rm foot}(A_-^pA_+^p)^{-1}$.  The overlay charges a product
intersection to the interval with the later left endpoint; each
interval is charged at most once to an interval whose left endpoint
precedes its own.  Its two sums are the selected total times the
opposite individual maximum and the unselected total times the
selected individual maximum.  This proves the displayed bound even
for arbitrarily short intervals.  It is not a product of physical
probabilities.  Counts of words or components are not multiplied
onto it afterwards.

For $k>j$ use exact reversal.  The selected backward half now has
\[
 h=k-j-1,\qquad n=j,\qquad l=k+1,
\]
and its bound is
\[
 C_{\rm foot}E_{\rm lc}
 \{K_{\rm loc}\sigma^{m-k-1}(\delta+\theta_{\rm r}^{j})
                             +H_{\rm lc}\sigma^k\}.
\]
At $k=j$ no separate selected-half assertion is needed: pay the
whole original boundary by
$C_{\rm foot}E_{\rm lc}H_{\rm lc}
 (\sigma^j+\sigma^{m-j-1})$.
Multiply these three bounds by $C_AT_{\rm lc}$, using
Lemma~\ref{lem:lc-turning}.  For $k\ne j$ the selected trace is
covered by the weak clearance intervals of Lemma~\ref{lem:lc-cover}.
Their newly included endpoints have zero source-angle measure by
its strict impact derivative; no separate positive trace is removed.
At $k=j$ the whole positive boundary was already retained.

Summing over every $k$ gives the sharper upper bound
\[
 \frac{K_{\rm bd}}{1-\sigma}
 \{2K_{\rm loc}\delta+
 (K_{\rm loc}+H_{\rm lc})
          (\theta_{\rm r}^j+\theta_{\rm r}^{m-1-j})\}.
\]
Here, for example, the unselected tail for $k<j$, together with
its $k=j$ endpoint, is bounded by
$H_{\rm lc}\sigma^{m-j-1}/(1-\sigma)$; the other tail is its
reversal.  Since $\sigma\le\theta_{\rm r}$, this proves
\eqref{eq:lc-old-boundary}.  The marked angular construction is a
one-copy parametrization after fixing the incident flight index and
mark; the complete two-sign source already includes every admissible
marked displacement and physical word.  Thus no unvisited disk,
collision word or shallow index has been omitted.

The new boundary is at $u_0=\delta$ on flight $j$.  The unselected
half-capacity bound and the same turning argument give
$K_{\rm bd}H_{\rm lc}(\sigma^j+\sigma^{m-j-1})$ for one incident
side.  We allow both sides explicitly, obtaining
\eqref{eq:lc-new-boundary}.  For an event with several qualifying
marks, assign each physical boundary germ to its first active mark
in a fixed ordering; discarding this assignment only enlarges the
complete positive marked sum.  A boundary where a selected foot
reaches a flight endpoint is impossible in this range: at an incident
contact its distance to a third center is at least $1-R$, whereas
$R+\delta_*<1-R$ uniformly.  The two incident disks themselves
are never counted as unvisited marks.  Thus the displayed list of
new physical boundaries is complete.
\end{proof}

\begin{lemma}[Selected endpoint-cap payment by the physical Jacobian]
\label{lem:lc-selected-cap}
The complete residual regular cap-side charge of
\eqref{eq:lc-selected-roof} is at most $C_{\rm cap,loc}\delta$.
\end{lemma}
\begin{proof}
On $p_0=a$, write $v=p_m$.  On every genuine regular root branch
of the endpoint map its inverse Jacobian obeys the exact identity
\[
                  u_x\,|dv|=|d\alpha_0|.
\]
The other component of the endpoint map is the unchanged coordinate
$p_0$, so this is simply its triangular determinant on the actual
fibre, with no probability normalization.  Moreover
$|d\arg(-a,v)|=a(a^2+v^2)^{-1}|dv|\le |dv|/a$.
By \eqref{eq:lc-boundary-density} the selected side charge is at
most $C_A/a$ times the selected $d\alpha_0$ mass.
For $j=0$ use \eqref{eq:lc-horizontal-initial}; for $j\ge1$ use
\eqref{eq:lc-source-mass} with $h=j-1$.  Each is at most
$M_{\rm cap}\delta$.

The side $p_0=-a$ has the same estimate.  For $p_m=\pm a$ apply
physical reversal, which changes the selected index to $m-1-j$;
the preceding initial and later-index cases cover it without an
additional factor.  Adding the four sides gives exactly the claimed
constant.  At simultaneous boundaries assign visited grazing first,
then remaining critical endpoints, then original unvisited constraints,
then the new level constraint, and finally the cap.  The residual
cap roots are therefore genuine regular roots where the displayed
Jacobian identity applies.  Corners and coordinate seams are assigned
once.  In particular this proof does not restrict an arbitrary
almost-everywhere version of a two-dimensional density to a null line.
\end{proof}

\subsection{The selected interior and the incoming critical endpoints}

\begin{theorem}[Complete physical selected roof height]
\label{thm:lc-selected-roof-height}
The measure \eqref{eq:lc-selected-roof} is absolutely continuous.
Its density $p_{m,R;j,\delta}$ satisfies, for every allowed
$m,R,j,\delta$,
\begin{equation}\label{eq:lc-roof-height}
 \|p_{m,R;j,\delta}\|_{L^\infty(\mathbb R)}
 \le C_{\rm clr,roof}
       \{\delta^{1/4}+\theta_{\rm r}^{j}
                             +\theta_{\rm r}^{m-1-j}\},
\end{equation}
where a fixed admissible choice is
\begin{equation}\label{eq:lc-roof-constant}
 C_{\rm clr,roof}
  =\frac{C_{\rm int}J_\times}{c_*R_-^{1/4}}
                         +C_{\rm side}+C_{\delta}+C_{\rm cap,loc}.
\end{equation}
The complete actual-word, physical mark and angular sums precede
every essential roof supremum and every supremum over the original
radius and exact labels.  All original joint return, displacement
and collision tests, the original section, and every Borel selector
$0\le H\le1$ may be retained by positive restriction.  The theorem
is for their unchanged physical time and makes no assertion about
an independent momentum process or a modified conditional event.
\end{theorem}
\begin{proof}
Absolute continuity follows by positive domination by the complete
cap measure in \eqref{eq:rbv-majorant} and its physical coarea.
For the localized interior variation use the exact pointwise bound
\eqref{eq:rbv-interior-pointwise}.  The full multiplier theorem
\eqref{eq:sm-full-functional} and the exact clearance flux give
\begin{align*}
 \int\sum_w\int_{\mathcal D_w(\theta)\cap E_{j,\delta,R}}
                 |\partial_t A_w|\,dt\,d\theta
 &\le\frac{C_{\rm int}}{c_*}
                   \int \one_{E_{j,\delta,R}}Z_m\,d\nu\\
 &\le\frac{C_{\rm int}J_\times}{c_*}
                            \nu(E_{j,\delta,R})^{1/4}\\
 &\le\frac{C_{\rm int}J_\times}{c_*R_-^{1/4}}\delta^{1/4}.
\end{align*}
Here $Z_m=I_m/|\nabla_qF|$ is the entire original singular
multiplier, and invariance together with
\eqref{eq:cf-exact-mean} gives
$\nu(E_{j,\delta,R})\le\int N_\delta\,d\nu=\delta/R$.
This calculation intersects actual observables on the same orbit;
it uses neither independence nor a positive incidence margin.
The factor $c_*^{-1}$ enters once, outside the unconditioned
localized multiplier estimate.

We give the endpoint-to-height step explicitly, since a small bulk
mass alone does not prove it.  Fix $m,R$.  In the complete analytic
physical endpoint charts of Lemma~\ref{lem:rbv-endpoints}, add the
finitely many signed impact and interior-foot conditions defining
$E_{j,\delta,R}$.  For almost every Gauss angle these give ordinary
open selected subintervals of the original physical roof intervals.
The same finite Boolean analytic description, compact chart cover
and one-dimensional slicing argument used there applies.  Equalities
have the explicit weak-boundary traces already counted above;
identically satisfied constraints are removed from the local
Boolean description.  Artificial seams and homogeneous cuts cancel
before absolute values are taken.  Distinct actual word sheets and
all genuine physical intervals remain in the sum.

The possible interval endpoints are thus visited grazing, a critical
point, an original unvisited tangency, a new clearance-$\delta$
level, or a cap side.  At a visited grazing the actual optical
factor $\omega$ tends to zero while $\chi\ge\lambda_\circ$;
hence $A=\omega/(4\pi c_*\chi)$ has zero physical coarea trace,
including a simultaneous critical limit.  This is the actual optical
trace assertion of Lemma~\ref{lem:rbv-endpoints}, not an inference
from the zero collision-area measure of the singular set.

Every remaining critical endpoint has positive visited incidences.
It is an incoming endpoint for the roof.  Indeed put $g=r e$ on
its fixed Gauss ray.  The positive reduced Hessian gives
\[
 \frac{dq}{dr}=G^{-1}e,\qquad
 \frac{dF}{dr}=r e^{\mathsf T}G^{-1}e>0\quad(r>0).
\]
At $r=0$ the roof is its lower one-sided limit.  Such a critical
point can therefore be a left endpoint of a selected interval, or
support no interval at all; it cannot be its right endpoint.  All
positive incoming critical fan values are retained in this statement.
No bound sets those values to zero, including a critical point on
an active clearance boundary.  Local charts and these conclusions
are used at fixed $m,R$ before nonnegative exhaustion; no uniform
Taylor radius is presumed.

On each selected fixed-angle interval $(a_q,b_q)$, use the actual
unweighted smooth density and its physical right trace.  Absolute
continuity on compact subintervals, followed by the proved finite
interior and outgoing boundary budgets, yields
\begin{equation}\label{eq:lc-right-endpoint}
 A_q(t)\le A_q(b_q^-)+\int_t^{b_q}|\partial_uA_q(u)|\,du,
                                      \qquad a_q<t<b_q.
\end{equation}
The whole physical roof lies in $[m\tau_-,m\tau_+]$, so there is
no unaccounted endpoint at infinity.  The right trace is either
zero visited grazing or a regular charge covered by
Lemmas~\ref{lem:lc-selected-boundary} and
\ref{lem:lc-selected-cap}; a remaining critical point was just
proved to be incoming.  For finitely many compact regular intervals
sum \eqref{eq:lc-right-endpoint}; integrate the Gauss angle and
then exhaust all actual intervals.  Tonelli and nonnegative
convergence apply to both upper budgets.  Their total bounds prove
\[
 \|p_{m,R;j,\delta}\|_\infty
 \le \frac{C_{\rm int}J_\times}{c_*R_-^{1/4}}\delta^{1/4}
       +C_{\rm side}(\delta+\theta_{\rm r}^j
                               +\theta_{\rm r}^{m-1-j})
       +C_\delta(\theta_{\rm r}^j+\theta_{\rm r}^{m-1-j})
       +C_{\rm cap,loc}\delta.
\]
Since $0<\delta\le\delta_*<1$, this implies
\eqref{eq:lc-roof-height}.  An arbitrary original Borel selector is
restored only after this positive bound and cannot increase its
density.  No variation of an indicator-cut density is asserted.
For each radius the conclusion is an essential supremum; uniform
constants then permit a supremum over radii.  A single exceptional
null set common to uncountably many radii is not required.
\end{proof}

The exponent $1/4$ comes from the actual localized singular multiplier,
whereas the two geometric tails come from the number of real collisions
between the marked flight and either endpoint.  The result controls the
selected density at a common physical roof.  It does not yet pay the
endpoint layers after multiplication by $m^2$, and it does not assert
abundance of the original nongrazing comparison pool.
